\documentclass[11pt,a4paper]{article}
\usepackage[T1]{fontenc}
\usepackage{lmodern,amsmath,amssymb,mathtools,booktabs,tabularx,microtype,xcolor}
\usepackage[bookmarks=false]{hyperref}
\usepackage[margin=24mm,headheight=14pt]{geometry}
\usepackage{fancyhdr}
\definecolor{ink}{HTML}{244B68}
\hypersetup{colorlinks=true,linkcolor=ink,pdftitle={POG DENSE: Expanding the Candidate Score}}
\pagestyle{fancy}\fancyhf{}\fancyhead[L]{\small POG\_DENSE: expanding the score}\fancyfoot[C]{\thepage}
\setlength{\parindent}{0pt}\setlength{\parskip}{5pt}\setlength{\emergencystretch}{2em}
\newcommand{\I}{\mathcal I}\newcommand{\pos}[1]{\left[#1\right]_+}
\newcommand{\off}{\operatorname{off}}\newcommand{\diag}{\operatorname{diag}}
\newcommand{\N}{\mathcal N}\newcommand{\sd}{\operatorname{sd}}
\newcommand{\why}{\par\textbf{Why this term?} }
\numberwithin{equation}{section}
\title{\textcolor{ink}{POG\_DENSE: Expanding the Candidate Score}\\[4pt]
\large From $\theta(c\mid S)$ to training bundles and item descriptions}
\author{}\date{}
\begin{document}
\begin{center}
{\LARGE\textcolor{ink}{POG\_DENSE: Expanding the Candidate Score}}\\[6pt]
{\large From $\theta(c\mid S)$ to training bundles and item descriptions}
\end{center}\medskip
Take a candidate item $c$ and a nonempty partially filled set $S$, with $c\notin S$.
Let $\I$ be the catalogue, $n=|\I|$, and $B\in\{0,1\}^{m\times n}$ the training table:
row $e$ is a past bundle, and $B_{ei}=1$ if it contains item $i$. Define the observed
frequency $f_i=\sum_e B_{ei}$. Our task is to compute $\theta(c\mid S)$ and rank candidates
by this score. This note follows the POG\_DENSE formulation with full item features.
We use the default uncapped known set $S$. The item categories and title, image, and
other description vectors are additional observed inputs; they cannot be recovered from
$B$ alone. Their construction is not assumed to follow from bundle membership.
Define $[x]_+=\max(x,0)$, $\varepsilon_8=10^{-8}$,
$\varepsilon_9=10^{-9}$, and $\varepsilon_{12}=10^{-12}$.
$\mathbf1[A]$ is one when statement $A$ is true and zero otherwise; $\lVert v\rVert_2$
is vector length, and $v\cdot w$ is the sum of products of corresponding entries.
The numerical constants below specify this version of the method. Their roles can be
explained, but their exact values are design choices rather than mathematical consequences.

\section{Start with the score: two explicitly weighted estimates}
Before giving the score, define $Y(c\mid S)$ as the combined evidence, $P(c\mid S)$ as
an estimate drawn mainly from categories and item descriptions, $s_0\ge0$ as its scale,
and $\kappa\ge0$ as the strength of the fallback. These are fixed after selection on
held-out examples. The final score is
\begin{equation}\label{theta}
\boxed{\theta(c\mid S)=
\begin{cases}
 Y(c\mid S),&\kappa=0,\\[2pt]
 \displaystyle\frac{f_c}{f_c+\kappa}Y(c\mid S)
 +\frac{\kappa}{f_c+\kappa}s_0P(c\mid S),&\kappa>0.
\end{cases}}
\end{equation}
\why The assumption behind this mixture is that more appearances make history-based evidence more dependable. Frequency alone does not prove that assumption for every item.
A candidate with no appearances receives $s_0P$ when the fallback is enabled. For example,
if $f_c=\kappa$, both estimates receive weight $1/2$. The result is a ranking score.
The next question is what goes into the two estimates $Y$ and $P$.

\subsection{Expand the combined evidence $Y$}
Let $p(c\mid S)$ denote a weighted sum of training-bundle matches. Let
$\delta(c\mid S)$ denote a second sum that includes only training bundles containing
at least two members of $S$. Let $L$ be a fixed item-to-item matrix built from $B$.
Section~\ref{bundledefs} gives their complete definitions. Here the important distinction
is the exact pair of raw calculations used for each source:
\begin{center}\small
\begin{tabular}{@{}lll@{}}\toprule
Source $h$ & Raw score $r_h^0$ & Raw score $r_h$ \\\midrule
bundle & $p_c$ & $p_c+\delta_c$ \\
propagated & $(Lp)_c$ & $(Lp)_c+(L\delta)_c$ \\
Each remaining source & $r_h$ & $r_h$ \\\bottomrule
\end{tabular}\end{center}
Here $p,\delta$ are column vectors, with $p_c=p(c\mid S)$ and
$\delta_c=\delta(c\mid S)$. Let $\mathcal H$ be the following ordered list:
\[
(\mathrm{cat,soft,text,image,bridge,bundle,propagated,pop,cooc,pmi,svd}).
\]
These are labels for the formulas expanded below, not additional assumptions. Let
$w_h\ge0$ be the selected weight of source $h$. Apply the same standardization rule
$\mathcal Z_S$ separately to each raw column, as defined in Section~\ref{standard}:
$z_h^0=\mathcal Z_S[r_h^0]$ and $z_h=\mathcal Z_S[r_h]$.
Let $\alpha\in[0,1]$ be a mixing coefficient, $\lambda_R\ge0$ an agreement coefficient, and $z_R$
the standardized pair-of-known-items score defined in Section~\ref{agreementdef}. Then
\begin{align}
Y_0(c\mid S)&=\sum_{h\in\mathcal H}w_h z_h^0(c\mid S),\\
Y_1(c\mid S)&=\sum_{h\in\mathcal H}w_h z_h(c\mid S),\\
Y(c\mid S)&=(1-\alpha)Y_0(c\mid S)+\alpha Y_1(c\mid S)+\lambda_R z_R(c\mid S).\label{Y}
\end{align}
\why The only difference between $Y_0$ and $Y_1$ is the explicit $\delta$ contribution
in the two rows of the table. This is a definition of this model: descriptions, category
scores and the other counts are held fixed during this comparison. It is not a claim
that only these two sources can describe relationships among several items. For instance,
$r_R$ also uses pairs of known items, but it is a separate term with its own coefficient $\lambda_R$.
Restricting $\alpha$ to $[0,1]$ makes $1-\alpha$ and $\alpha$ nonnegative weights adding to one; values outside this interval would extrapolate rather than average. Requiring $\lambda_R\ge0$ keeps a positive relative agreement score from being penalized through a negative coefficient. These are modeling restrictions. At $\alpha=0$, the two rows use their left-hand formulas; at $\alpha=1$, they use their right-hand
formulas. Intermediate $\alpha$ mixes the separately standardized scores. It does not replace
$p$ by $p+\alpha\delta$ before standardization.

\subsection{Expand the fallback $P$}
Define $\mathcal Q=\{\mathrm{cat,soft,text,image,bridge}\}$ restricted to available
sources, and let $\mathcal A=\{h\in\mathcal H:w_h>0\}$ after removal of rejected weights.
The prior-source set is defined precisely by
\[
\mathcal P=\begin{cases}\mathcal Q\cap\mathcal A,&\mathcal Q\cap\mathcal A\ne\varnothing,\\
\mathcal Q,&\mathcal Q\cap\mathcal A=\varnothing.\end{cases}
\]
This is a prescribed choice of prior sources, not a set forced by a theorem. For $h\in\mathcal P$, let $\widetilde w_h=w_h$ when its
positive weight was retained, and $\widetilde w_h=0.5$ otherwise. An omitted source has $w_h=0$ in $Y$; the value $0.5$ in $P$ is a separate default rule. Let $z_{\rm pop}$ denote
the standardized popularity score. Then
\begin{equation}\label{P}
 \boxed{P(c\mid S)=\sum_{h\in\mathcal P}\widetilde w_hz_h(c\mid S)
       +\lambda_R z_R(c\mid S)+w_{\rm pop}z_{\rm pop}(c\mid S).}
\end{equation}
\why For an item with $f_c=0$, its row in $B^\top B$ and in the associated training-only
item matrices is zero. Such rows cannot provide item-specific positive shared-history
support. Categories and supplied feature vectors can still distinguish that item, so
$\mathcal Q$ selects those kinds of sources. This is the motivation for this restriction,
not a proof that this particular restriction is best. Bridge and category tables still
use training history; the fallback is not independent of $B$. Selected popularity and agreement also enter this estimate. The scale
$s_0$ makes its overall standard deviation comparable to that of $Y$. We now expand the standardized
scores $z_h$ used in both equations.

\section{Expand $z_h$: put each source on a common scale}\label{standard}
Let $r_h(i\mid S)$ be the raw score of source $h$ for catalogue item $i$. Define
$\varepsilon_9=10^{-9}$, the catalogue mean $\mu_h$, and sample standard deviation $\sigma_h$ by
\begin{align}
 \mu_h(S)&=\frac1n\sum_{i\in\I}r_h(i\mid S),\\
 \sigma_h(S)&=\sqrt{\frac1{n-1}\sum_{i\in\I}(r_h(i\mid S)-\mu_h(S))^2},\\
 \boxed{z_h(c\mid S)}&\boxed{=\frac{r_h(c\mid S)-\mu_h(S)}{\sigma_h(S)+\varepsilon_9}.}\label{z}
\end{align}
Denote this column operation by $\mathcal Z_S[r_h]$. Apply it separately to $r_h^0$ and to the agreement score $r_R$.
The sums include the known items; they are excluded only when producing suggestions.
\why A shared-history count and an image similarity have different units. This conversion
expresses each as above or below its own catalogue average. A constant raw score becomes
zero and carries no distinction between candidates. In particular, normalizing $p+\delta$
is different from separately normalizing $p$ and $\delta$ and adding them. The inclusion of known items and the divisor $n-1$ are fixed conventions of this formulation; neither is required by the definition of a ranking task.

The remaining task is to define each raw source in the ordered list $\mathcal H$ already given.
We expand the bundle and propagated scores first, since they also define the connections used
by the other parts of the score.

\section{Expand the two training-bundle formulas}\label{bundledefs}
\subsection{Direct bundle matches: $p(c\mid S)$ and the extra overlap $\delta(c\mid S)$}
Define bundle size $d_e=\sum_i B_{ei}$, known-item weight
$u_s=(f_s+5)^{-1/2}$, and bundle weight $v_e=(d_e+1)^{-1/2}$. Let $o_e$ be the weighted
overlap with the known items, and $k_e$ their unweighted overlap count:
\begin{align}
 o_e(S)&=v_e\sum_{s\in S}B_{es}u_s,& k_e(S)&=\sum_{s\in S}B_{es},\\
 p(c\mid S)&=\sum_e B_{ec}o_e(S),&
 \delta(c\mid S)&=\sum_e B_{ec}o_e(S)\,1.5[k_e(S)-1]_+ .\label{delta}
\end{align}
Here $[x]_+=\max(x,0)$. The two raw versions entering $Y_0$ and $Y_1$ are
\begin{equation}\label{bundle}
 \boxed{r_{\rm bundle}^0(c\mid S)=p(c\mid S),\qquad
 r_{\rm bundle}(c\mid S)=p(c\mid S)+\delta(c\mid S).}
\end{equation}
\why $p$ adds evidence from past bundles containing the candidate and known items.
For example, the factor $(f_s+5)^{-1/2}$ decreases when a known item's count increases,
and $(d_e+1)^{-1/2}$ decreases as a bundle becomes larger. The offsets 5 and 1 prevent
singular weights and soften the adjustment. The exponent $1/2$ and these offsets are
fixed choices; many alternative decreasing functions would also be mathematically valid. More explicitly, a bundle with $k_e=0$ or $1$ contributes zero to $\delta$; $k_e=2$
contributes $1.5 B_{ec}o_e$; $k_e=3$ contributes $3 B_{ec}o_e$. This condition concerns
how many members of $S$ occur in a particular training row, as specified by $k_e(S)$.
The candidate must also occur in that row, through $B_{ec}$. Thus $\delta=0$ for $|S|=1$,
and also whenever no training row contains two known items. The factor $1.5$ and the linear
increase in $k_e-1$ specify a chosen strength and shape; they are not derived optimal values.

\subsection{Spread history: apply the table $L$ to the bundle scores}
Let $D=\diag((f_i+5)^{-1/2})$ and $E=\diag((d_e+1)^{-1})$ be diagonal tables, and let
$\off(A)$ mean $A$ with its diagonal replaced by zero. Define
\begin{equation}\label{L}
 L=\off(DB^\top EBD).
\end{equation}
For a vector $q$ with one entry per item, $(Lq)_c=\sum_i L_{ci}q_i$.
Write $p$ and $\delta$ for the vectors of the scores just defined. Then
\begin{equation}
 \boxed{r_{\rm propagated}^0(c\mid S)=(Lp)_c,\qquad
 r_{\rm propagated}(c\mid S)=(Lp)_c+(L\delta)_c.}
\end{equation}
\why A candidate can benefit from its connection to another item that already has strong
support. $L$ passes evidence through shared training bundles while excluding direct
self-transfer. Entry by entry, for $c\ne i$,
\[
L_{ci}=\frac{1}{\sqrt{f_c+5}\sqrt{f_i+5}}\sum_e\frac{B_{ec}B_{ei}}{d_e+1},\qquad L_{cc}=0.
\]
Thus $(Lp)_c$ adds $p_i$ from each other item $i$, weighted by an explicitly defined
shared-bundle connection. This is what ``propagated'' means here. A previously unseen
item has a zero row in $L$ and receives no such transfer.

\section{Expand the other history sources}
\subsection{Popularity, shared counts, and unexpectedly frequent pairs}
Define $C=\off(B^\top B)$, $g_i=\sum_j C_{ij}+\varepsilon_9$, and
$Z_C=\sum_{i,j}C_{ij}$. Introduce $\varepsilon_{12}=10^{-12}$ and define
\begin{equation}
 M_{ij}=\begin{cases}
 [\log(C_{ij}Z_C/(g_i g_j)+\varepsilon_{12})]_+,&C_{ij}>0,\\
 0,&C_{ij}=0.
 \end{cases}
\end{equation}
The raw scores are
\begin{equation}
 \boxed{r_{\rm pop}(c\mid S)=\log(1+f_c),\quad
 r_{\rm cooc}(c\mid S)=\sum_{s\in S}C_{cs},\quad
 r_{\rm pmi}(c\mid S)=\sum_{s\in S}M_{cs}.}
\end{equation}
\why Popularity gives a general frequency preference. Shared counts measure direct
pair evidence. $M$ instead asks whether the pair is strong relative to the overall shared
counts of its two items. Taking a logarithm compresses large ratios; taking the positive part discards below-expected associations from this source rather than assigning them a negative raw penalty. These are distinct views of the same training table.

\subsection{A short summary of the pair table}
Set $r=\min(128,n-2)$. The limit 128 is a chosen summary width, not an intrinsic rank of the data. Let $V_r$ contain the leading $r$ right singular vectors of
$M$: orthogonal directions that, for an exact decomposition, give its best rank-$r$ approximation in squared error. A numerical approximation to these directions is sufficient to define the fitted feature table.
Introduce $\varepsilon_8=10^{-8}$ and the row-normalization rule
$\N(F)_i=F_i/(\lVert F_i\rVert_2+\varepsilon_8)$. Define $U=\N(MV_r)$; if $M$ is
zero, use $U=0$. The score is
\begin{equation}
 \boxed{r_{\rm svd}(c\mid S)=U_c\cdot\left(\frac1{|S|}\sum_{s\in S}U_s\right).}
\end{equation}
\why This summarizes many shared-history relationships in a shorter vector per item.
The dot product compares the candidate's summary with the mean summary of the known items.
It supplements the full pair table rather than replacing all history scores with a short vector.

\section{Expand the text, image, and bridged-description sources}
Let $F^{t},F^{i},F^{d}$ denote the available title, image, and description feature tables,
with rows normalized by $\N$. Let $\ell(i)$ be item $i$'s clothing category. For modality
$\nu\in\{t,i,d\}$, define the category average $\bar F^\nu_a$ and the remaining feature $D^\nu$:
\begin{equation}
 \bar F^\nu_a=\frac{\sum_{i:\ell(i)=a}F^\nu_i}{|\{i:\ell(i)=a\}|},\qquad
 D^\nu_i=\frac{F^\nu_i-\bar F^\nu_{\ell(i)}}
 {\lVert F^\nu_i-\bar F^\nu_{\ell(i)}\rVert_2+\varepsilon_8}.
\end{equation}
\why Removing the category average focuses on what distinguishes items within a clothing
category. This subtracts the mean feature vector of the category. It does not guarantee statistical independence from category, nor does it prove that the remaining features measure compatibility.

Let $D_* =\N([D^t|D^i|D^d])$, where brackets mean concatenating feature columns.
The table $LD_*$ is a weighted sum of other items' feature rows, using exactly the entries of $L$ above. Define
\begin{equation}
 F_b=\N\bigl(D_*+0.5\N(LD_*)\bigr).
\end{equation}
The three selectable feature scores are
\begin{align}
 r_{\rm text}(c\mid S)&=D^t_c\cdot\frac1{|S|}\sum_{s\in S}D^t_s,\\
 r_{\rm image}(c\mid S)&=D^i_c\cdot\frac1{|S|}\sum_{s\in S}D^i_s,\\
 \boxed{r_{\rm bridge}(c\mid S)}&\boxed{=F_{b,c}\cdot\frac1{|S|}\sum_{s\in S}F_{b,s}.}
\end{align}
\why Text and images provide information even for items with no training appearances.
The bridge adds information from historically connected items to each item's own features.
Description features contribute to this bridge. Their absence as a separately weighted source is a restriction of the chosen POG\_DENSE model, not a consequence of their being mathematically redundant. The bridge coefficient $0.5$ is a fixed choice. Unavailable modalities and their separately weighted sources are omitted. If title features are unavailable, the agreement term is omitted and $\lambda_R=0$.

\section{Expand the category sources}
\subsection{Named clothing categories}
Let $A_{ea}=\sum_i B_{ei}\mathbf1[\ell(i)=a]$ count category $a$ in bundle $e$.
Define $H=\off(A^\top A)$, $h_a=\sum_e A_{ea}+\varepsilon_9$, and $H_+=\sum_a h_a$.
Define the category-pair table
\begin{equation}
 Q_{ab}=\begin{cases}[\log(H_{ab}H_+/(h_a h_b)+\varepsilon_{12})]_+,&H_{ab}>0,\\0,&H_{ab}=0.\end{cases}
\end{equation}
Let $n_a^{(1)}$ and $n_a^{(2)}$ be the numbers of training bundles containing at least
one and at least two items of category $a$, respectively. Define
$\rho_a=(n_a^{(2)}+0.5)/(n_a^{(1)}+1)$. Then
\begin{equation}
 \boxed{r_{\rm cat}(c\mid S)=\sum_{s\in S}Q_{\ell(s),\ell(c)}
 -0.5\sum_{s\in S}\mathbf1[\ell(s)=\ell(c)]\mathbf1[\rho_{\ell(s)}<0.15](1-\rho_{\ell(s)}).}
\end{equation}
\why The first part rewards category combinations seen together more often than their
category counts suggest. The second discourages repeating a category that rarely appears
twice in a training bundle. The counts come from training bundles; the smoothing constants $0.5,1$, the penalty multiplier $0.5$, and threshold $0.15$ are prescribed choices. The ratio $\rho_a$ estimates how frequently the category is repeated among bundles containing it; smoothing keeps it defined even for a category absent from training.

\subsection{Groups suggested by item features}
Named categories are coarse. Let $X=\N([F^t|F^i|F^d])$ and choose
$k=\max(2,\min(400,n-1))$ feature groups using alternating assignment to a nearest center and center updates. This aims to reduce within-group squared distances, without requiring a global minimum. The limit 400 is a chosen number of groups, not the number of clothing categories.
Let $v_a$ be the center of feature group $a$, normalized with $\N$. The constant $\tau_s=0.05$ controls how sharply membership favors the most similar center: smaller values make membership more concentrated. It is a chosen setting. Define
item $i$'s degree of membership in feature group $a$ by
\begin{equation}
 t_{ia}=X_i\cdot v_a/\tau_s,\qquad
 J_{ia}=\frac{\exp(t_{ia}-\max_b t_{ib})}
 {\sum_b\exp(t_{ib}-\max_b t_{ib})+\varepsilon_9}.
\end{equation}
Let $T=BJ$ be the training-bundle group totals. Define group-pair counts, catalogue group
masses, and their total by
\begin{equation}
 H^s=\pos{\off\left(T^\top T-\sum_e\sum_{i:B_{ei}=1}J_i^\top J_i\right)},\qquad
 a_u=\sum_iJ_{iu},\qquad A_s=\sum_u a_u.
\end{equation}
Finally, define $Q^s_{uv}=[\log(H^s_{uv}A_s/(a_u a_v+\varepsilon_9)+\varepsilon_{12})]_+$.
The raw score is
\begin{equation}
 \boxed{r_{\rm soft}(c\mid S)=(JQ^s)_c\cdot\left(\sum_{s\in S}J_s\right).}
\end{equation}
\why An item can partly belong to several feature groups. Their training-bundle
relationships may capture combinations missed by a single clothing label. In fact, before the diagonal is removed, the table being counted is
\[
T^\top T-\sum_e\sum_i B_{ei}J_i^\top J_i
=\sum_e\sum_{i\ne j}B_{ei}B_{ej}J_i^\top J_j.
\]
This is a sum over distinct items in the same training bundle. The subtraction therefore
removes an item's contribution paired with itself; otherwise a feature group
could appear to complement itself merely because it describes the same item twice.

\section{Expand the pair-of-known-items term $r_R$}\label{agreementdef}
The terms in $\mathcal H$ are now defined. The remaining term in $Y$ and $P$ is agreement
among known items. Define $V=\diag((d_e+1)^{-1/2})$ and
$W=\off(B^\top VBD)$. For known item $s$, define its support for the candidate by
\begin{equation}
 q_{cs}=W_{cs}+0.15[D^t_c\cdot D^t_s]_+,\qquad A_c=\sum_{s\in S}q_{cs}.
\end{equation}
For two distinct known items define their similarity $G_{st}$ and their pair weight $K_{st}$:
\begin{equation}
 G_{st}=0.5\mathbf1[\ell(s)=\ell(t)]+0.5[D^t_s\cdot D^t_t]_+,\qquad
 K_{st}=[1-G_{st}]_+,\qquad K_{ss}=0.
\end{equation}
The raw agreement score, standardized using equation~\eqref{z}, is
\begin{equation}
 \boxed{r_R(c\mid S)=\frac{\sum_{s\in S}\sum_{t\in S}q_{cs}K_{st}q_{ct}}{A_c+0.05}.}
\end{equation}
\why This rewards a candidate supported by several known items, while reducing the extra
credit for support from very similar known items. Both orders of a pair enter the sum.
For $S=\{s,t\}$ with $s\ne t$, the numerator is exactly
$2q_{cs}q_{ct}K_{st}$. It is zero if either known item supplies zero support, or if $K_{st}=0$.
The denominator controls the scale of this product. For a single known item it is zero
because $K_{ss}=0$. It is not a logical requirement to match all of $S$: for a larger
set, any supported pair can contribute. The values $0.15$, $0.5$, and $0.05$ are fixed
choices specifying the feature contribution, similarity mixture, and denominator smoothing.
No derivation requires those particular values. The table $W$ has zero diagonal; the bundle score $p$ defined
earlier retains self-contributions before final candidate exclusion. These differ when
computing the catalogue-wide scaling and propagated scores.

\section{Where the fixed settings come from}
The expansion of $\theta$ is now complete except for how $w_h,\alpha,\lambda_R,s_0,\kappa$ and
$\mathcal P$ are chosen. They are selected on separate examples with known hidden items,
then held fixed when scoring new requests. The hidden items in the final test examples
are not used for these choices. In particular, $\alpha$ could mathematically be any value in
$[0,1]$, and $\lambda_R$ any nonnegative value. The finite sets below are the search candidates
chosen for this method, so that fitting uses finitely many evaluations. They are not the
only admissible values, and no claim of optimal spacing or exhaustive optimization follows.
The same distinction applies to the weight grid, tolerances, four-pass limit and clipping
bounds. Their numerical values specify this procedure; a mathematical justification for
these exact values is not supplied by the scoring equations.

\textbf{Weights and mixture.} Visit the sources in the order
cat, soft, text, image, bridge, bundle, propagated, pop, cooc, pmi, svd. Starting at zero,
try each weight in $\{0,0.25,0.5,1,1.5\}$, skipping the all-zero weight vector and using $Y_0$ and keeping changes that improve
mean NDCG@20 by more than $2\times10^{-4}$, for at most four passes. For request $j$, let $G_j$ be its nonempty hidden-item set, let $\pi_j(r)$ be its
$r$th ranked non-known item, and let $K_j=\min(20,n-|S_j|)$. Define
\[
\mathrm{NDCG}_{20,j}=\frac{\sum_{r=1}^{K_j}\mathbf1[\pi_j(r)\in G_j]/\log_2(r+1)}
{\sum_{r=1}^{\min(20,|G_j|)}1/\log_2(r+1)}.
\]
Mean NDCG@20 is its average across held-out requests.
Try removing bridge and soft-category sources if doing so loses at most $3\times10^{-4}$, retaining at least one positive-weight source.
Choose $\mathcal P$ from positive-weight category/soft-category/text/image/bridge sources,
falling back to all available ones if necessary. Then compare $\alpha\in\{0,0.5,1\}$ and
$\lambda_R\in\{0,0.2,0.35\}$, retaining improvements exceeding $5\times10^{-5}$.
\why The method learns which evidence is useful for the completion task without assuming
that every source deserves equal importance.

\textbf{Scale.} Across all catalogue items and held-out requests, let $\sd$ denote sample
standard deviation. With the selected weights and mixture, first estimate
\begin{equation}
 s_* =\frac{\sd(Y)+\varepsilon_9}{\sd(P)+\varepsilon_9}.
\end{equation}
\why This aligns the overall size of variations in the two estimates before mixing them.

\textbf{History adjustment.} Let there be $v$ held-out requests indexed by $j$. For each
item with $f_i>0$, define $x_i=1/f_i$ and
$e_i=v^{-1}\sum_j(Y(i\mid S_j)-s_*P(i\mid S_j))^2$.
Let $\bar x,\bar e$ be their means across such items. Fit a straight line to these values:
\begin{align}
 b_*&=\frac{\sum_i(x_i-\bar x)(e_i-\bar e)}{\sum_i(x_i-\bar x)^2},&
 a_*&=\bar e-b_*\bar x,\\
 \kappa_*&=\operatorname{clip}\left(\frac{\max(b_*,10^{-4})}{\max(a_*,10^{-4})},0.05,50\right).
\end{align}
If fewer than two items have history, or the denominator is at most $10^{-12}$, use only
$\kappa=0$. Otherwise compare $\kappa\in\{0,0.05,0.5,1,\kappa_*,5\}$ with
$s_0\in\{s_*,1\}$, retaining NDCG@20 improvements greater than $5\times10^{-5}$.
\why Let $\theta_i^*$ be an unobserved suitability, and $\mu_i=s_*P(i\mid S)$.
Assume $\theta_i^*=\mu_i+\xi_i$ and $Y_i=\theta_i^*+\eta_i$, with independent,
zero-mean normal errors of variances $\sigma_0^2$ and $\sigma_e^2/f_i$, respectively.
For $f_i>0$ these assumptions imply
\[
\mathbb E[(Y_i-\mu_i)^2\mid f_i]=\sigma_0^2+\frac{\sigma_e^2}{f_i},
\qquad
\mathbb E[\theta_i^*\mid Y_i,\mu_i,f_i]
=\frac{f_iY_i+\kappa\mu_i}{f_i+\kappa},\quad
\kappa=\frac{\sigma_e^2}{\sigma_0^2}.
\]
This explains both the fitted line and the rational weights in equation~\eqref{theta}.
The assumptions are a model of score errors, not a consequence of the observed counts.
The slope and intercept estimate the two variances; positivity bounds and clipping turn
their ratio into one search candidate. Held-out ranking quality chooses the final setting.
The squared residuals compare $Y$ and $s_*P$, not either score with a binary item label.

\section{Read the expansion back upward}
For a candidate $c$ and known set $S$, compute the required raw components from $B$,
category labels, and item features. Apply equation~\eqref{z} to each component, substitute
the results into $Y$ and $P$ in equations~\eqref{Y}--\eqref{P}, then substitute those into
$\theta$ in equation~\eqref{theta}. Exclude all items already in $S$ and rank the others by
descending $\theta$, resolving exact ties by increasing item ID. The highest $K$ scores
are the requested completion. Sources with zero weight and no role in the fallback
contribute nothing to this calculation.
\end{document}
